\section{Introduction}
As trains become automated, the need for operators to monitor and control them remotely is growing. To do this safely, the operator needs a live video feed from the train that is stable and fast enough to react to what is happening on the track. This thesis looks at how such a video streaming system can be built, tested and evaluated, with a focus on finding configurations that work well for remote train operation.

\subsection{Background and Motivation}
As the railway sector transitions toward higher Grades of Automation (GoA), particularly GoA 3 and GoA 4, remote driving becomes an increasingly critical operational capability. Remote driving is not only a natural consequence of reduced onboard personnel at higher GoA levels, but serves as a mandatory safety layer within the ATO architecture enabling operators to intervene when automated systems encounter situations they cannot resolve independently, such as unexpected track conditions or degraded operational modes~\cite{orr2024, europesrail2025_d3}.

A fundamental requirement for safe remote driving is the availability of reliable, low-latency video streaming. When the operator is physically separated from the train, the video feed becomes the primary source of situational awareness, and its quality and responsiveness directly determine whether timely and safe intervention is possible~\cite{mejias2024, myhrvold2025}. Industry experience further demonstrates that remote operation can support a range of practical tasks including shunting, depot movements, and vehicle preparation provided that the communication link meets sufficient performance criteria~\cite{fp2r2dato2023d5_4}.

Understanding the relationship between latency and human perception is equally important in this context. Research on human factors in remote operation indicates that video delay, bitrate, and stream stability all influence an operator's ability to detect obstacles, interpret dynamic scenes, and execute corrective actions under time pressure~\cite{brandenburger2023, neumeier2019}. Identifying the latency boundaries where safe remote operation remains usable is therefore a key step toward enabling higher automation levels in railway systems~\cite{myhrvold2025}.

This thesis is developed within an ongoing research collaboration at NTNU, in partnership with VTI and other European railway organizations, contributing to the technical groundwork required to make remote train driving a safe and practical reality in future railway deployments.

\subsection{Problem}

The current remote train operation setup relies on a commercial product, which imposes significant limitations on the project's ability to test and develop the system further. As a closed external solution, it restricts access to the underlying pipeline, making it difficult to adjust configurations, measure internal latency, or experiment with different streaming approaches. In a research context where iterative testing and configuration control are essential, this represents a fundamental obstacle. In addition, licensing such a product over the course of a long term research project represents a considerable and ongoing cost.

At the same time, there is limited existing research that addresses video streaming specifically in the context of remote train operation. Most work on low latency video transmission is either too general or focused on other domains, leaving a gap when it comes to understanding what configurations and compromises are most suitable for railway use cases~\cite{myhrvold2025}.

Building a custom streaming solution also comes with its own set of challenges. Network conditions such as signal strength, interference, and varying load introduce unpredictable variability in transmission quality. Choosing the right codec, transport protocol, and encoder settings requires systematic testing, as each parameter affects latency, quality, and resource usage in ways that are not always predictable. Balancing video quality against latency is a difficult comparison and testing these differences in settings on real hardware in a realistic setup adds further complexity. Without a flexible and open pipeline, it is difficult to isolate and measure where latency is introduced and how different configurations affect overall system performance.

\subsection{Objectives}

The main objective of this thesis is to design and implement a low-latency video streaming pipeline tailored specifically for remote train operation, serving as a practical replacement for the current commercial solution. The pipeline should be functional, flexible, and open enough to support continued development and testing within the project.
More specifically, the objectives are to:

\begin{enumerate}
    \item O1: Develop a working GStreamer-based streaming pipeline that meets the low-latency requirements of remote train operation
    \item O2: Implement a latency measurement framework that allows end-to-end delay to be captured, logged, and compared across different configurations
    \item O3: Evaluate the pipeline across multiple configurations, including different codecs, transport protocols and network settings, to identify the interplay between latency and video quality
    \item O4: Design the system to be dynamic and extensible, allowing additional cameras and new configurations to be added with minimal effort
    \item O5: Deliver a stable and usable prototype that is good enough for practical use during the current development period and serves as a foundation for future work
\end{enumerate}

The result should be a system that gives the project full control over the streaming pipeline, removes dependency on external licensed software, and produces reliable data to guide future configuration decisions.

\subsection{Research Questions}
The following research questions guide the experimental work of this thesis:
\begin{itemize}
    \item RQ1: What end-to-end latency and video quality can be achieved in a GStreamer-based video streaming pipeline for remote train operation, and which codec and transport configuration performs best?
    \item RQ2: How do network conditions, including transport protocol and available bandwidth, affect streaming latency and reliability in a remote train operation context?
    \item RQ3: What are the practical interplay between latency, video quality, and resource usage across different streaming configurations, and which configuration is most suitable for remote train operation?
\end{itemize}

\subsection{Scope and Limitations}

The scope of this thesis covers the design, implementation and experimental evaluation of a custom video streaming pipeline for remote train operation. The system is developed as a proof of concept and is not intended as a deployment ready solution. Development and pre testing was conducted in a controlled indoor environment, while final testing was performed from a moving vehicle to simulate realistic train movement conditions.

The experimental evaluation focuses on the parameters identified as most relevant to remote train operation based on prior research \cite{myhrvold2025}, and selected based on time and hardware availability. The configurations tested include codec selection (H.264, H.265 and MJPEG), transport protocol (UDP and TCP), encryption overhead (SRTP), bandwidth conditions, resolution, and GPU versus CPU based encoding. While GStreamer is open for a large number of additional parameters to be adjusted, expanding further was not feasible within the time constraints of a single master thesis.

Latency is measured end to end between the point where an RTP packet enters the sender pipeline and the point where the corresponding decoded frame arrives at the receiver display sink. Camera capture latency and screen rendering latency are not included in these measurements, as they occur outside the instrumented part of the pipeline and are considered largely constant across configurations. Clock synchronization between the sender and receiver relies on internet based time sources, introducing a small margin of uncertainty estimated at a few milliseconds, which is acknowledged when interpreting absolute latency values.

The thesis does not include testing on an actual operational train. A moving vehicle was used to simulate train conditions, but real world railway infrastructure, operational speeds, and trackside network coverage are outside the scope of this work and represent a natural next step for future validation.

The thesis does not include formal safety certification or a RAMS analysis in accordance with EN 50126. While the system is developed with safety critical use in mind, evaluating compliance with railway safety standards is considered outside the scope of this work and is left for future development phases.

\begin{comment}
With RAMS:
The system is developed with reference to EN 50126, the European railway standard for reliability, availability, maintainability and safety. While a full RAMS analysis is beyond the scope of a single master thesis, relevant safety considerations from the standard are used to inform design decisions and contextualize the results within a safety critical operational environment.
\end{comment}

\subsection{Outline of Structure}

The remainder of this thesis is organized as follows:

Chapter 2 presents the theoretical background, covering remote train operation and the grades of automation, the role of video and the latency requirements that apply to it, the fundamentals of video streaming and compression including the H.264, H.265 and MJPEG codecs, the transport protocols used, the individual sources of latency across a streaming pipeline, clock synchronisation, the GStreamer framework, the relevant network technologies, and the performance metrics used to evaluate the system. 

Chapter 3 describes the design of the system itself, setting out the requirements, the iterative development methodology, the overall architecture and components, the GStreamer sender and receiver pipelines, and the measurement design used to capture each metric without disturbing the live stream. 

Chapter 4 explains the methodology, including the research and literature approach, the reasoning behind transcoding, the design of the experiments and the full test matrix, the equipment and procedure, the controls applied to limit hidden variables, and the reliability and validity of the work. 

Chapter 5 presents the experimental results for the transport protocol tests and the codec configuration tests across all five metrics, with interpretation deliberately held back. 

Chapter 6 then discusses those results, interpreting them through the three research questions, comparing them against the 3GPP latency target, examining the performance of the transport protocols and codec configurations, drawing out the interplay between latency, quality and resource usage, arriving at a configuration recommendation, and reflecting on the limitations and the relation to prior work. 

Chapter 7 concludes by summarising the main findings and outlining directions for future work.

